\Needspace{7\baselineskip} % REV03: subsection follows the complete electronic argument.
\subsection{Transport lorentzien et représentation électromagnétique}

La section électronique fournit une structure et une coordonnée
avant son application à une interaction. La réalisation lorentzienne
de l'incidence cubique détermine, avec la convention temporelle déclarée,
l'espace des formes dans lequel cette interaction peut s'exprimer.
Pour la formuler, on fixe une réalisation locale orientée \((\mathcal M,g)\)
de dimension quatre, avec des fibres tangentes identifiées au modèle
\(Q_\square\). On spécifie en outre une connexion \(\nabla\) avec
\(\nabla g=0\), une un-forme potentielle \(A\in\Omega^1(\mathcal M)\),
sa courbure \(F=dA\) et une représentation de charge de la section.
Ces données sont les arguments de la représentation physique qui suit.

La dualité de Hodge satisfait \(\star^2=-I\) sur les deux-formes. La séparation électrique et magnétique correspond à un choix temporel dans la décomposition \(1+3\) ; \(F\) et \(\star F\) conservent leur relation au sein du même espace de deux-formes. La neutralité d'une représentation de charge signifie \(q=0\), tandis que la section peut conserver orientation, phase et transport interne. L'absence de couplage électrique dans cette représentation a un contenu distinct de l'absence d'information de phase.

Dans cette réalisation dimensionnelle, soient \(m_e>0\), \(q\) la charge,
\(\gamma(s)\) une trajectoire paramétrée et \(u=\dot\gamma\)
sa quadrivitesse. La force de Lorentz s'exprime par
\begin{equation}
 m_e\nabla_u u^\mu=qF^\mu{}_{\nu}u^\nu.
 \label{eq:lorentz-electron}
\end{equation}
L'antisymétrie de \(F\) implique
\(F_{\mu\nu}u^\mu u^\nu=0\). Par conséquent,
\[
 \frac{d}{ds}g(u,u)=2g(\nabla_u u,u)=0.
\]
La normalisation de la trajectoire est conservée dès lors que \(\nabla\)
est métrique et que \eqref{eq:lorentz-electron} est satisfaite. La masse utilisée
est la réalisation dimensionnelle de la section électronique précédente ;
le paramètre de couplage est identifié dans la convention électromagnétique
adoptée. La déduction concrète du potentiel et d'une configuration
de champ constitue un problème supplémentaire assorti de conditions de frontière
propres. Le résultat exposé ici détermine le domaine, les données
de la représentation et la conservation qu'implique son équation.

L'inversion des routes du TPK et l'orientation d'une ligne électronique
offrent alors une comparaison précise : chacune possède son involution,
son transport et sa représentation. L'antécédent de Wheeler et Feynman
situe historiquement la question ; l'équivalence de deux réalisations
se décide au moyen de leurs applications et conditions de compatibilité.
Cette disposition permet de maintenir le lien entre l'origine arithmétique de l'électron
et la formulation de son interaction, en explicitant la fonction
de chaque étape.
